\documentclass{pnpcards}

\EnableCardDebugFrames

\typeout{pnpcards geometry: paper=\the\paperwidth\space x \the\paperheight}

\newcommand{\GeometryLine}[1]{%
    \rule{0pt}{8mm}\texttt{#1}\par
}

\begin{document}

\RunCard{A6 geometry}{%
    \GeometryLine{FRONT CONTENT LINE 01}
    \GeometryLine{FRONT CONTENT LINE 02}
    \GeometryLine{FRONT CONTENT LINE 03}
    \GeometryLine{FRONT CONTENT LINE 04}
    \GeometryLine{FRONT CONTENT LINE 05}
    \GeometryLine{FRONT CONTENT LINE 06}
    \GeometryLine{FRONT CONTENT LINE 07}
    \GeometryLine{FRONT CONTENT LINE 08}
    \GeometryLine{FRONT CONTENT LINE 09}
    \GeometryLine{CONTINUATION LINE 10}
    \GeometryLine{CONTINUATION LINE 11}
    \GeometryLine{CONTINUATION LINE 12}
}

\end{document}

% Expected Output
% The PDF has two A6 landscape pages (148 mm x 105 mm), with all four outer
% frame edges 5 mm from the page edges.
% Page 1: The class-generated title frame measures 138 mm x 15 mm. The front
%         content frame measures 138 mm x 78 mm, with a 2 mm gap below the
%         title frame. Together the frames and gap use the full 95 mm between
%         the 5 mm top and bottom margins.
%         FRONT CONTENT LINE 01 onward appears in the front content frame.
% Page 2: The class-generated continuation frame measures 138 mm x 95 mm,
%         fills the full printable height, has no title frame, and starts
%         with CONTINUATION LINE 10. It is visibly taller than the front
%         content frame. No frame crosses the 5 mm margins.
